% =====================================================================
%  Documentation technique — AcadémiePlus
%  Plateforme adaptative d'apprentissage des mathématiques
%  Compilation : pdfLaTeX (Overleaf) — encodage UTF-8
% =====================================================================
\documentclass[11pt,a4paper]{report}

% ---------- Encodage & langue ----------
\usepackage[utf8]{inputenc}
\usepackage[T1]{fontenc}
\usepackage[french]{babel}
\usepackage{lmodern}
\usepackage{microtype}

% ---------- Mise en page ----------
\usepackage[a4paper,margin=2.4cm,headheight=15pt]{geometry}
\usepackage{graphicx}
\usepackage{xcolor}
\usepackage{fancyhdr}
\usepackage{titlesec}
\usepackage{enumitem}
\usepackage{booktabs}
\usepackage{longtable}
\usepackage{array}
\usepackage{tabularx}
\usepackage{multirow}
\usepackage{amsmath,amssymb}
\usepackage{hyperref}
\usepackage{tcolorbox}
\usepackage{listings}
\usepackage{fontawesome5}
\tcbuselibrary{skins,breakable}

% ---------- Palette (charte AcadémiePlus) ----------
\definecolor{apblue}{HTML}{1E3A8A}     % bleu profond
\definecolor{apaccent}{HTML}{2563EB}   % bleu accent
\definecolor{apgreen}{HTML}{059669}    % vert réussite
\definecolor{apgray}{HTML}{475569}     % gris texte
\definecolor{aplight}{HTML}{EFF4FF}    % fond clair
\definecolor{apdark}{HTML}{0F172A}     % fond foncé
\definecolor{apcode}{HTML}{F1F5F9}     % fond code

% ---------- Hyperliens ----------
\hypersetup{
  colorlinks=true,
  linkcolor=apaccent,
  urlcolor=apaccent,
  citecolor=apaccent,
  pdftitle={Documentation technique - AcademiePlus},
  pdfauthor={Equipe AcademiePlus}
}

% ---------- En-têtes / pieds ----------
\pagestyle{fancy}
\fancyhf{}
\renewcommand{\headrulewidth}{0.4pt}
\renewcommand{\footrulewidth}{0pt}
\fancyhead[L]{\small\textcolor{apblue}{\textbf{AcadémiePlus}}}
\fancyhead[R]{\small\textcolor{apgray}{\leftmark}}
\fancyfoot[C]{\small\textcolor{apgray}{\thepage}}

% ---------- Style des titres ----------
\titleformat{\chapter}[display]
  {\normalfont\huge\bfseries\color{apblue}}
  {\filright\Large\color{apaccent}Chapitre~\thechapter}{8pt}{\Huge}
\titlespacing*{\chapter}{0pt}{-10pt}{24pt}

\titleformat{\section}
  {\normalfont\Large\bfseries\color{apblue}}{\thesection}{1em}{}
\titleformat{\subsection}
  {\normalfont\large\bfseries\color{apaccent}}{\thesubsection}{1em}{}

% ---------- Encadrés ----------
\newtcolorbox{infobox}[1][]{
  enhanced, breakable, colback=aplight, colframe=apaccent,
  boxrule=0.8pt, arc=3pt, left=10pt, right=10pt, top=8pt, bottom=8pt,
  fonttitle=\bfseries\color{white},
  coltitle=white, attach boxed title to top left={yshift=-2mm,xshift=4mm},
  boxed title style={colback=apaccent,arc=2pt}, title={#1}}

\newtcolorbox{notebox}[1][\faInfoCircle\ Note]{
  enhanced, breakable, colback=white, colframe=apgreen,
  boxrule=0.8pt, arc=3pt, left=10pt, right=10pt, top=8pt, bottom=8pt,
  fonttitle=\bfseries\color{white}, coltitle=white,
  attach boxed title to top left={yshift=-2mm,xshift=4mm},
  boxed title style={colback=apgreen,arc=2pt}, title={#1}}

% ---------- Style code ----------
\lstdefinestyle{ap}{
  backgroundcolor=\color{apcode},
  basicstyle=\ttfamily\footnotesize\color{apdark},
  keywordstyle=\color{apaccent}\bfseries,
  commentstyle=\color{apgreen}\itshape,
  stringstyle=\color{apgreen},
  numberstyle=\tiny\color{apgray},
  numbers=left, numbersep=8pt, frame=single,
  rulecolor=\color{apgray!40},
  breaklines=true, showstringspaces=false, tabsize=2,
  captionpos=b, xleftmargin=14pt, framexleftmargin=10pt}
\lstset{style=ap}

\newcommand{\filepath}[1]{\textcolor{apgray}{\ttfamily\small #1}}

% =====================================================================
\begin{document}

% ---------- Page de titre ----------
\begin{titlepage}
\centering
\vspace*{1.5cm}
{\color{apaccent}\rule{\linewidth}{2pt}}\\[0.6cm]
{\Huge\bfseries\color{apblue} AcadémiePlus}\\[0.4cm]
{\LARGE\color{apgray} Plateforme adaptative d'apprentissage\\ des mathématiques}\\[0.5cm]
{\color{apaccent}\rule{\linewidth}{2pt}}\\[1.4cm]

{\Large\bfseries Documentation technique complète}\\[0.4cm]
{\large\color{apgray} Architecture, fonctionnalités, base de données et moteur adaptatif}\\[2.2cm]

\begin{tcolorbox}[width=0.8\linewidth,colback=aplight,colframe=apblue,arc=4pt]
\centering
\begin{tabular}{rl}
\textbf{Projet} & AcadémiePlus (\texttt{academieplus.lovable.app})\\
\textbf{Type} & Application web SPA (React + Supabase)\\
\textbf{Domaine} & Soutien scolaire en mathématiques\\
\textbf{Public} & Élèves, Parents, Pédagogues, Administrateurs\\
\textbf{Version doc.} & 1.0\\
\textbf{Date} & Juin 2026\\
\end{tabular}
\end{tcolorbox}

\vfill
{\large\color{apgray}\faGraduationCap\quad Document destiné aux futures équipes de développement}\\[0.4cm]
{\small\color{apgray} Ce document décrit l'intégralité des choix techniques, des fonctionnalités\\ et de la structure de données du projet.}
\end{titlepage}

% ---------- Sommaire ----------
\tableofcontents
\thispagestyle{empty}
\clearpage

% ---------- Chapitres ----------
\chapter{Introduction}

\section{Présentation du projet}
\textbf{AcadémiePlus} est une plateforme web de soutien scolaire spécialisée en
\textbf{mathématiques}, destinée aux élèves du système éducatif algérien
(du primaire jusqu'au baccalauréat). L'application combine du contenu de cours
structuré, des exercices et quiz interactifs, un \textbf{moteur d'apprentissage
adaptatif} qui ajuste la difficulté au niveau réel de l'élève, et un
\textbf{assistant pédagogique IA} (chatbot) encadré par le programme scolaire.

\begin{infobox}[\faBullseye\ Objectif pédagogique]
Offrir à chaque élève un parcours \emph{personnalisé} : le système mesure en
permanence le niveau de l'élève (sur 100), détecte ses lacunes, génère des
exercices calibrés et fournit des commentaires et recommandations automatiques.
Les parents suivent la progression de leurs enfants via un tableau de bord dédié.
\end{infobox}

\section{Publics cibles et rôles}
La plateforme distingue quatre rôles applicatifs :
\begin{itemize}[leftmargin=1.4cm]
  \item[\faUserGraduate] \textbf{Élève (\texttt{student})} --- suit les cours, passe le test
        de positionnement, réalise exercices/quiz, dialogue avec l'IA.
  \item[\faUsers] \textbf{Parent (\texttt{parent})} --- relie ses enfants à son compte et
        consulte leurs rapports de progression.
  \item[\faChalkboardTeacher] \textbf{Pédagogue (\texttt{pedago})} --- rédige et gère le
        contenu des cours, exercices, quiz et examens (espace éditorial).
  \item[\faUserShield] \textbf{Administrateur (\texttt{admin})} --- supervise les
        utilisateurs, les abonnements, les codes d'activation et l'ensemble du contenu.
\end{itemize}

\section{Niveaux scolaires couverts}
Le contenu est organisé par niveau scolaire (énumération \texttt{school\_level}) :
\begin{center}
\begin{tabular}{ll}
\toprule
\textbf{Code interne} & \textbf{Niveau} \\
\midrule
\texttt{5eme\_primaire} & 5\textsuperscript{e} année primaire \\
\texttt{1ere\_cem} & 1\textsuperscript{re} année CEM (collège) \\
\texttt{2eme\_cem} & 2\textsuperscript{e} année CEM \\
\texttt{3eme\_cem} & 3\textsuperscript{e} année CEM \\
\texttt{4eme\_cem} & 4\textsuperscript{e} année CEM (BEM) \\
\texttt{premiere} & 1\textsuperscript{re} année secondaire \\
\texttt{seconde} & 2\textsuperscript{e} année secondaire \\
\texttt{terminale} & 3\textsuperscript{e} année secondaire (BAC) \\
\bottomrule
\end{tabular}
\end{center}

Chaque niveau secondaire peut être subdivisé en \textbf{filières} (sciences,
lettres, gestion, math-techniques\dots), ce qui détermine les chapitres affichés.

\section{Comment lire cette documentation}
\begin{itemize}
  \item Le \textbf{chapitre 2} décrit l'architecture globale.
  \item Le \textbf{chapitre 3} détaille la pile technique et les versions.
  \item Le \textbf{chapitre 4} explique l'authentification et la gestion des rôles.
  \item Le \textbf{chapitre 5} liste les fonctionnalités (pages) écran par écran.
  \item Le \textbf{chapitre 6} décrit le moteur adaptatif (calcul du niveau).
  \item Les \textbf{chapitres 7 et 8} documentent la base de données et le MCD.
  \item Le \textbf{chapitre 9} décrit les fonctions serveur (Edge Functions IA).
  \item Les \textbf{chapitres 10 et 11} traitent la sécurité et le déploiement.
\end{itemize}

\chapter{Architecture générale}

\section{Vue d'ensemble}
AcadémiePlus est une \textbf{application monopage (SPA)} React, hébergée et
servie en statique, qui communique avec un \textbf{backend Supabase}
(PostgreSQL + Auth + Storage + Edge Functions). Aucune partie serveur «~maison~»
n'est maintenue~: toute la logique métier vit soit dans le client React, soit
dans des \textbf{fonctions Edge Deno} déployées sur l'infrastructure Supabase.

\begin{verbatim}
        +-------------------------------------------------------------+
        |                     NAVIGATEUR (Client)                     |
        |                                                             |
        |   React 18 + Vite + TypeScript + Tailwind + shadcn/ui      |
        |   - Pages & composants (src/pages, src/components)          |
        |   - Hooks métier (src/hooks)                                |
        |   - Services d'accès données (src/services)                 |
        |   - Moteur adaptatif pur (src/lib/levelEngine.ts)          |
        +-----------------------------+-------------------------------+
                                      |  HTTPS (supabase-js)
                                      v
        +-------------------------------------------------------------+
        |                   SUPABASE (Backend cloud)                  |
        |                                                             |
        |   Auth (JWT)   PostgreSQL (RLS)   Storage (avatars)         |
        |                                                             |
        |   Edge Functions (Deno) :                                   |
        |     - lovable-chat (assistant IA)                           |
        |     - generate-adaptive-content                             |
        |     - generate-lesson-comment / remediation / report ...    |
        +-----------------------------+-------------------------------+
                                      |  HTTPS
                                      v
        +-------------------------------------------------------------+
        |          FOURNISSEURS IA (via clés stockées en secret)      |
        |   Lovable AI Gateway (Gemini) | Groq (Llama) | Gemini API   |
        +-------------------------------------------------------------+
\end{verbatim}

\section{Principes architecturaux}
\begin{itemize}
  \item \textbf{Séparation des responsabilités} : les composants d'UI
        (\filepath{src/components}) ne contiennent pas la logique de calcul.
        Le calcul du niveau est isolé dans des \emph{fonctions pures}
        (\filepath{src/lib/levelEngine.ts}), testables unitairement.
  \item \textbf{Accès données centralisé} : les appels Supabase passent par des
        \emph{services} (\filepath{src/services}) et des \emph{hooks}
        (\filepath{src/hooks}).
  \item \textbf{Sécurité côté base} : chaque table est protégée par des règles
        \textbf{RLS} (Row-Level Security) ; les rôles sont stockés dans une table
        dédiée \texttt{user\_roles} (jamais sur le profil).
  \item \textbf{IA déportée} : les appels aux modèles d'IA se font exclusivement
        depuis des Edge Functions, afin de ne jamais exposer les clés au client.
\end{itemize}

\section{Organisation des dossiers}
\begin{longtable}{p{5cm}p{9.5cm}}
\toprule
\textbf{Dossier} & \textbf{Contenu} \\
\midrule
\endhead
\filepath{src/pages/} & Les pages routées (une par écran). \\
\filepath{src/pages/editorial/} & Espace de rédaction pour pédagogues/admins. \\
\filepath{src/components/} & Composants réutilisables (UI, course, dashboard\dots). \\
\filepath{src/components/ui/} & Bibliothèque shadcn/ui (boutons, dialogues\dots). \\
\filepath{src/components/course/} & Cœur pédagogique : leçons, exercices, quiz. \\
\filepath{src/hooks/} & Hooks React métier (adaptatif, profil, abonnement\dots). \\
\filepath{src/lib/} & Logique pure : moteur de niveau, helpers, validation. \\
\filepath{src/services/} & Accès aux données Supabase encapsulés. \\
\filepath{src/contexts/} & Contextes React (authentification). \\
\filepath{src/data/} & Données statiques des programmes (chapitres par niveau). \\
\filepath{src/i18n/} & Internationalisation (français / arabe). \\
\filepath{src/integrations/supabase/} & Client Supabase et types générés. \\
\filepath{supabase/functions/} & Fonctions Edge (Deno) côté serveur. \\
\filepath{supabase/migrations/} & Migrations SQL versionnées de la base. \\
\bottomrule
\end{longtable}

\chapter{Pile technique et versions}

\section{Technologies principales}
L'application repose sur une pile \textbf{React/Vite/TypeScript} moderne avec
\textbf{Tailwind CSS} pour le style et \textbf{Supabase} pour le backend.

\begin{longtable}{p{4.2cm}p{2.2cm}p{8cm}}
\toprule
\textbf{Technologie} & \textbf{Version} & \textbf{Rôle} \\
\midrule
\endhead
React & 18.3 & Bibliothèque d'interface (composants, hooks). \\
Vite & 5.4 & Bundler / serveur de développement rapide. \\
TypeScript & 5.8 & Typage statique de tout le code. \\
Tailwind CSS & 3.4 & Système de design utilitaire (tokens sémantiques). \\
shadcn/ui (Radix) & --- & Composants accessibles (dialogues, menus, onglets\dots). \\
React Router DOM & 6.30 & Routage côté client (SPA). \\
TanStack React Query & 5.83 & Cache et synchronisation des requêtes serveur. \\
Supabase JS & 2.97 & Client base de données / Auth / Storage / Functions. \\
\bottomrule
\end{longtable}

\section{Bibliothèques notables}
\begin{longtable}{p{4.5cm}p{10cm}}
\toprule
\textbf{Paquet} & \textbf{Usage dans le projet} \\
\midrule
\endhead
\texttt{i18next} / \texttt{react-i18next} & Internationalisation FR/AR + détection de langue. \\
\texttt{katex}, \texttt{react-katex}, \texttt{rehype-katex}, \texttt{remark-math} & Rendu des formules mathématiques (\LaTeX) dans les leçons et le chat. \\
\texttt{react-markdown}, \texttt{remark-gfm}, \texttt{rehype-raw} & Rendu du contenu Markdown (cours, solutions, commentaires IA). \\
\texttt{@tiptap/*} & Éditeur de texte riche WYSIWYG pour l'espace éditorial. \\
\texttt{recharts} & Graphiques de progression (tableaux de bord, analytics). \\
\texttt{framer-motion} & Animations (compagnon-bot, transitions). \\
\texttt{jspdf}, \texttt{html2canvas} & Export PDF (cours, rapports parents). \\
\texttt{react-hook-form} + \texttt{zod} & Formulaires et validation de schémas. \\
\texttt{@dnd-kit/*} & Glisser-déposer (réordonnancement de sections éditoriales). \\
\texttt{date-fns} & Manipulation des dates. \\
\texttt{lucide-react}, \texttt{react-icons} & Jeux d'icônes. \\
\texttt{sonner}, \texttt{vaul}, \texttt{embla-carousel} & Toasts, tiroirs, carrousels. \\
\bottomrule
\end{longtable}

\section{Outils de développement et qualité}
\begin{itemize}
  \item \textbf{ESLint 9} + \texttt{typescript-eslint} --- qualité et cohérence du code.
  \item \textbf{Vitest 3} + \texttt{@testing-library/react} --- tests unitaires/composants.
  \item \textbf{Playwright 1.57} --- tests de bout en bout (navigateur).
  \item \textbf{PostCSS / Autoprefixer} --- chaîne CSS.
  \item \textbf{@vitejs/plugin-react-swc} --- compilation rapide via SWC.
\end{itemize}

\section{Scripts npm}
\begin{lstlisting}[language=bash,caption={package.json --- scripts}]
npm run dev        # serveur de developpement Vite
npm run build      # build de production
npm run build:dev  # build en mode developpement
npm run lint       # analyse ESLint
npm run test       # tests Vitest (une passe)
npm run test:watch # tests en mode surveillance
\end{lstlisting}

\begin{notebox}
Le projet est construit avec l'écosystème \textbf{Lovable}. Les fonctions Edge
sont déployées automatiquement à chaque modification ; le fichier \filepath{.env}
est auto-rempli avec \texttt{VITE\_SUPABASE\_URL},
\texttt{VITE\_SUPABASE\_PUBLISHABLE\_KEY} et \texttt{VITE\_SUPABASE\_PROJECT\_ID}.
\end{notebox}

\chapter{Authentification et gestion des rôles}

\section{Système d'authentification}
L'authentification s'appuie sur \textbf{Supabase Auth} (jetons JWT). Le contexte
React \filepath{src/contexts/AuthContext.tsx} centralise l'état de session~:
\begin{itemize}
  \item Au chargement, \texttt{supabase.auth.getSession()} récupère la session
        courante ; un abonnement \texttt{onAuthStateChange} maintient l'état à jour.
  \item Il expose \texttt{user}, \texttt{session}, \texttt{loading},
        \texttt{signOut()}, \texttt{hasRole(role)} et \texttt{isAdmin()}.
\end{itemize}

\section{Redirection intelligente après connexion}
Après chaque connexion, le contexte oriente l'utilisateur selon son rôle et son
état~:
\begin{itemize}
  \item \textbf{Sans rôle} $\rightarrow$ \filepath{/complete-profile} (compléter le profil).
  \item \textbf{Pédago / Admin} $\rightarrow$ \filepath{/liste-cours}.
  \item \textbf{Parent} $\rightarrow$ \filepath{/parent-dashboard}.
  \item \textbf{Élève sans test de positionnement} $\rightarrow$
        \filepath{/learning-assessment} (test initial obligatoire).
  \item \textbf{Élève} (cas général) $\rightarrow$ \filepath{/cours/math}.
\end{itemize}

\begin{infobox}[\faKey\ Détection du test de positionnement]
La fonction \texttt{hasCompletedPlacementAssessment} vérifie l'existence d'une
ligne \texttt{student\_scores} «~globale~» (\texttt{lesson\_id IS NULL}). Si
elle est absente, l'élève est forcé de passer le test initial avant d'accéder
aux cours. Une compatibilité ascendante évite de re-forcer les anciens comptes.
\end{infobox}

\section{Stockage sécurisé des rôles}
Conformément aux bonnes pratiques de sécurité, les rôles \textbf{ne sont jamais
stockés sur le profil}. Ils résident dans une table dédiée \texttt{user\_roles}
et sont vérifiés côté base via une fonction \texttt{SECURITY DEFINER}~:

\begin{lstlisting}[language=SQL,caption={Vérification de rôle sans récursion RLS}]
CREATE OR REPLACE FUNCTION public.has_role(_user_id uuid, _role app_role)
RETURNS boolean LANGUAGE sql STABLE SECURITY DEFINER
SET search_path TO 'public' AS $$
  SELECT EXISTS (
    SELECT 1 FROM public.user_roles
    WHERE user_id = _user_id AND role = _role
  )
$$;
\end{lstlisting}

L'énumération des rôles est \texttt{app\_role = \{student, parent, admin, pedago\}}.

\section{Création automatique de profil}
Un trigger \texttt{handle\_new\_user} s'exécute à l'inscription~: il crée la ligne
\texttt{profiles} (avec normalisation du niveau scolaire et des métadonnées) et
insère le rôle dans \texttt{user\_roles} à partir des métadonnées d'inscription.

\section{Protection des routes}
Le composant \filepath{src/components/ProtectedRoute.tsx} encapsule chaque route
privée. Il accepte~:
\begin{itemize}
  \item \texttt{allowedRoles} — liste blanche de rôles autorisés ;
  \item \texttt{requireAdmin} — accès réservé aux administrateurs ;
  \item \texttt{blockAdmin} — interdit l'accès aux administrateurs (ex.\ page Compte).
\end{itemize}

\section{Liaison parent–enfant}
Un parent relie un enfant grâce à un \textbf{code de liaison} (\texttt{linking\_code}
dans \texttt{profiles}). La relation est stockée dans \texttt{parent\_child\_links}
avec un statut (\texttt{pending}, \texttt{active}, \texttt{rejected}). La fonction
\texttt{is\_parent\_of} sécurise l'accès des parents aux données de leurs enfants.

\chapter{Fonctionnalités (écran par écran)}

L'application est une SPA dont chaque écran correspond à une route déclarée dans
\filepath{src/App.tsx}. Les routes sont protégées par rôle via
\texttt{ProtectedRoute}. Cette section décrit chaque page par domaine.

\section{Pages publiques}
\begin{longtable}{p{3.2cm}p{3.4cm}p{7.6cm}}
\toprule \textbf{Page} & \textbf{Route} & \textbf{Description} \\ \midrule \endhead
Accueil & \texttt{/} & Page marketing (hero, atouts, tarifs, FAQ, CTA). \\
Authentification & \texttt{/auth} & Connexion / inscription (e-mail + mot de passe, choix du rôle, localisation, SSO). \\
Contact & \texttt{/contact} & Formulaire de contact. \\
Mentions légales & \texttt{/mentions-legales} & Page légale statique. \\
Confidentialité & \texttt{/politique-confidentialite} & Politique de confidentialité. \\
404 & \texttt{*} & Page «~non trouvée~». \\
\bottomrule
\end{longtable}

\section{Onboarding}
\begin{longtable}{p{3.2cm}p{3.4cm}p{7.6cm}}
\toprule \textbf{Page} & \textbf{Route} & \textbf{Description} \\ \midrule \endhead
Compléter profil & \texttt{/complete-profile} & Complétion post-SSO (rôle, niveau, filière, localisation, date de naissance). \\
Test de positionnement & \texttt{/learning-assessment} & \textbf{Test initial} de l'élève : 5 questions adaptées au niveau scolaire ; le score initialise le niveau (ligne globale \texttt{student\_scores}). \\
Onboarding Bot & \texttt{/bot-onboarding} & Questionnaire conversationnel du style d'apprentissage (visuel / textuel / pratique) animé par un bot bilingue. \\
\bottomrule
\end{longtable}

\section{Espace élève}
\begin{longtable}{p{3.2cm}p{3.2cm}p{7.8cm}}
\toprule \textbf{Page} & \textbf{Route} & \textbf{Description} \\ \midrule \endhead
Tableau de bord & \texttt{/dashboard} & Synthèse de l'élève : niveau, activité récente, commentaires IA, recommandations. \\
Liste des cours & \texttt{/liste-cours} & Catalogue des matières / chapitres selon le niveau. \\
Cours & \texttt{/cours/:subjectId} & \textbf{Interface centrale} : arbre chapitres/leçons, contenu adaptatif, chatbot flottant, quiz et exercices, export PDF. \\
Révision & \texttt{/revision/:subjectId} & Mode fiches (flashcards) de révision, cartes retournables. \\
Simulation & \texttt{/simulation/:subjectId} & Examen blanc chronométré (30 min), score et bilan par chapitre. \\
Examens (trimestre) & \texttt{/exams} & Sélection du trimestre. \\
Liste d'examens & \texttt{/exams/list} & Sujets d'examen par trimestre, mode chronométré. \\
Remédiation & \texttt{/remediation} & Activités ciblées sur les lacunes (3 exercices + 2 quiz via IA). \\
Compte & \texttt{/account} & Gestion du compte : avatar, abonnement (jours restants, pause), mot de passe. \\
Abonnements & \texttt{/abonnements} & Offres et état de l'abonnement, résiliation. \\
Paiement & \texttt{/paiement} & Tunnel de paiement. \\
Parrainage & \texttt{/parrainage} & Programme de parrainage (code, suivi, récompenses). \\
Profil & \texttt{/profile} & Vue profil (avatar, infos scolaires, statistiques). \\
Factures & \texttt{/factures} & Historique des factures. \\
Mes informations & \texttt{/mes-informations} & Édition des informations personnelles. \\
Mes données & \texttt{/mes-donnees-personnelles} & Consultation / suppression des données (RGPD). \\
\bottomrule
\end{longtable}

\section{Espace parent}
\begin{longtable}{p{3.2cm}p{3.6cm}p{7.4cm}}
\toprule \textbf{Page} & \textbf{Route} & \textbf{Description} \\ \midrule \endhead
Tableau parent & \texttt{/parent-dashboard} & Liste des enfants liés, progression par matière, rapports IA, ajout d'enfant par code, export PDF. \\
Cours (vue parent) & \texttt{/parent-cours/:childId} & Vue en lecture seule de la progression d'un enfant. \\
\bottomrule
\end{longtable}

\section{Espace administration}
\begin{longtable}{p{3.4cm}p{3.6cm}p{7.2cm}}
\toprule \textbf{Page} & \textbf{Route} & \textbf{Description} \\ \midrule \endhead
Admin & \texttt{/admin} & Gestion des utilisateurs (recherche, activation, suppression), statistiques, audit. \\
Abonnements admin & \texttt{/admin/abonnements} & Gestion des abonnements et des offres. \\
Analytics & \texttt{/analytics} & Journal d'activité, comptages, répartition des actions. \\
FAQ Admin & \texttt{/faq-admin} & Édition des FAQ de la page d'accueil. \\
Génération de contenu & \texttt{/content-generation} & Génération en masse d'exercices/quiz (progression en flux). \\
Éditeur de leçon & \texttt{/lecon/:lessonId} & Édition riche d'une leçon (TipTap, enrichissement IA, versions). \\
\bottomrule
\end{longtable}

\section{Espace éditorial (pédagogues / admins)}
\begin{longtable}{p{3.6cm}p{3.8cm}p{6.8cm}}
\toprule \textbf{Page} & \textbf{Route} & \textbf{Description} \\ \midrule \endhead
Tableau éditorial & \texttt{/editorial} & Liste des chapitres/filières avec statut et date de modification. \\
Éditeur de cours & \texttt{/editorial/cours/:id} & Édition complète : sections, texte riche, formules, détection de conflit. \\
Prévisualisation & \texttt{.../preview} & Aperçu d'une leçon telle que vue par l'élève. \\
Historique & \texttt{.../historique} & Historique des versions, restauration. \\
Comparaison & \texttt{.../compare} & Comparaison de deux versions (diff). \\
Révision & \texttt{/editorial/revision} & File de relecture / validation. \\
Médiathèque & \texttt{/editorial/mediatheque} & Bibliothèque de médias (Storage). \\
Gestion équipe & \texttt{/editorial/equipe} & Gestion des membres de l'équipe pédagogique. \\
\bottomrule
\end{longtable}

\section{Composants pédagogiques clés}
\begin{itemize}
  \item \filepath{AdaptiveLesson} — affiche la leçon avec bascule \emph{vidéo / texte}
        selon le style d'apprentissage ; rendu KaTeX des formules.
  \item \filepath{LessonActivityTabs} — organise les activités en trois étapes
        pédagogiques~: \<اكتشف> (découvrir), \<افهم> (comprendre), \<تعمّق>
        (approfondir, contenu IA). L'accès à l'approfondissement est conditionné
        à l'abonnement ou à 3 bonnes réponses.
  \item \filepath{ChapterMathExercises} — exercices à réponse libre, clavier
        mathématique, indices, validation serveur, verrou 3\,s sur erreur,
        détection d'hésitation.
  \item \filepath{ChapterMathQuiz} — quiz à choix multiples avec confirmation,
        suivi du temps et bilan final.
  \item \filepath{ChatBot} — assistant IA flottant et déplaçable (texte, image,
        saisie vocale FR/AR), encadré par le programme.
\end{itemize}

\section{Assistant IA — chatbot}
Le chatbot (\filepath{src/components/ChatBot.tsx}) appelle la fonction Edge
\texttt{lovable-chat}. Il est~:
\begin{itemize}
  \item \textbf{contextualisé} : reçoit le contenu du chapitre courant et la liste
        des leçons (pour des liens de navigation) ;
  \item \textbf{cadré} : répond uniquement en mathématiques, dans le programme ;
  \item \textbf{multimodal} : accepte texte, images et saisie vocale ;
  \item \textbf{évaluateur} : peut poser des exercices ; un marqueur caché
        \texttt{[[EVAL:correct|difficulté|concept]]} permet d'alimenter
        directement le moteur de niveau (\filepath{src/lib/chatEvalTracker.ts}).
\end{itemize}

\subsection{Quotas d'usage}
Le hook \filepath{src/hooks/useChatLimits.ts} applique des quotas pour les comptes
gratuits~: \textbf{10 messages} et \textbf{3 images} par jour (table
\texttt{chat\_usage}). Les abonnés actifs ont un accès illimité. L'historique est
sauvegardé (avec anti-rebond de 2\,s) dans \texttt{chat\_conversations}, en
retirant les données d'image pour alléger le stockage.

\section{Suivi parental et rapports}
Les parents lient leurs enfants via un code, puis consultent des
\textbf{rapports de progression} générés par IA (fonction
\texttt{generate-parent-report}). Une tâche planifiée
(\texttt{scheduled-parent-reports}) produit ces rapports périodiquement et
envoie des alertes en cas d'inactivité prolongée de l'enfant.

\chapter{Moteur d'apprentissage adaptatif}
\label{chap:adaptatif}

Ce chapitre répond directement à la question \emph{«~comment est calculé le
niveau de l'élève~?~»}. Toute la logique de calcul est concentrée dans des
fonctions pures de \filepath{src/lib/levelEngine.ts}, ce qui la rend testable
et indépendante de l'interface.

\section{Principe général}
Chaque élève possède un \textbf{niveau} compris entre \textbf{5 et 100} pour
chaque leçon (et un niveau global par chapitre). Le niveau évolue après chaque
réponse selon un modèle inspiré du système \textbf{ELO}, pondéré par la
difficulté, le temps de réponse et des \textbf{signaux comportementaux}.

\section{Niveau initial = test de positionnement}
À la création de son compte, l'élève passe un \textbf{test diagnostic}
(\filepath{src/pages/LearningAssessment.tsx}). Son score initial est enregistré
dans une ligne globale de \texttt{student\_scores} (\texttt{lesson\_id = NULL}).

\begin{infobox}[\faInfoCircle\ Pas de 50 par défaut]
Le niveau initial \textbf{n'est plus une valeur arbitraire de 50}. La fonction
\texttt{getPlacementLevel(userId)} (\filepath{src/lib/initialLevel.ts}) lit le
\texttt{current\_level} de la ligne de positionnement et l'utilise pour amorcer
chaque nouvelle leçon. À défaut (ancien compte), la valeur 50 sert de repli.
\end{infobox}

\section{Étape 1 --- Ajustement ELO par réponse}
La fonction \texttt{computeDelta} calcule la variation $\Delta$ du niveau après
chaque réponse.

\subsection{Pondérations}
\begin{center}
\begin{tabular}{ccc}
\toprule
\textbf{Difficulté} & \textbf{Poids} & \textbf{Temps médian attendu (s)} \\
\midrule
1 & 0,6 & 30 \\
2 & 0,8 & 45 \\
3 & 1,0 & 60 \\
4 & 1,3 & 90 \\
5 & 1,6 & 120 \\
\bottomrule
\end{tabular}
\end{center}

\subsection{Score attendu et delta de base}
Le «~score attendu~» suit une logistique de type ELO~:
\[
E = \frac{1}{1 + 10^{\,(20 \cdot d - L)/40}}
\]
où $d$ est la difficulté et $L$ le niveau courant. Puis~:
\[
\Delta =
\begin{cases}
\mathrm{clamp}\big(\mathrm{round}(3(1-E)\cdot w),\,1,\,5\big) & \text{si correct}\\[4pt]
\mathrm{clamp}\big(\mathrm{round}(-2E\cdot w),\,-5,\,-1\big) & \text{si faux}
\end{cases}
\]
avec $w$ le poids de difficulté.

\subsection{Modulation par le temps de réponse}
\begin{itemize}
  \item Réponse correcte \emph{très rapide} ($< 0{,}6 \times$ médian) : $\Delta \times 1{,}2$.
  \item Réponse correcte \emph{très lente} ($> 2 \times$ médian) : $\Delta \times 0{,}8$.
  \item Réponse fausse après un temps très long ($> 120$\,s) : $\Delta \times 1{,}3$.
\end{itemize}

\section{Étape 2 --- Signaux comportementaux}
Le moteur enrichit le delta en tenant compte de la manière dont l'élève a
répondu~:

\begin{longtable}{p{5.5cm}p{8.5cm}}
\toprule
\textbf{Signal} & \textbf{Effet sur le gain/la perte} \\
\midrule
\endhead
Tentatives multiples (réussite après erreurs) & gain $\times\,0{,}85^{(\text{tentatives}-1)}$ (plancher $+1$). \\
Indice \textbf{préventif} (consulté avant toute réponse) & gain $\times\,0{,}70$ ($-30\%$). \\
Indice \textbf{curatif} (consulté après une erreur) & gain $\times\,0{,}40$ ($-60\%$). \\
Abandon avec indice curatif & perte $\times\,1{,}3$. \\
\bottomrule
\end{longtable}

Le delta final est re-borné à $[+1, +6]$ (réussite) ou $[-7, -1]$ (échec).

\subsection{Hésitation sans réponse soumise}
La fonction \texttt{hesitationAdjustment(kind, difficulty)} gère deux cas~:
\begin{itemize}
  \item \textbf{«~timeout~»} (exercice ouvert $>1$\,min sans réponse puis quitté)~:
        $-1$, ou $-2$ si difficulté $\geq 4$.
  \item \textbf{«~erase~»} (saisie puis effacement total puis sortie)~: $0$ sur le
        niveau (enregistré uniquement dans le profil pour affiner le suivi).
\end{itemize}
Ces drapeaux sont cumulés dans le champ JSON \texttt{assessment\_data}
(\texttt{hesitation\_timeout}, \texttt{hesitation\_erase}).

\section{Étape 3 --- Score composite (calibration de l'IA)}
Pour calibrer la difficulté du contenu généré, on calcule un \textbf{score
composite} (\texttt{computeComposite})~:
\[
C = 0{,}40 \times \text{taux pondéré} + 0{,}35 \times L + 0{,}25 \times \text{taux quiz}
\]
borné à $[5, 100]$. Ce score sert de \texttt{difficulty\_level} transmis aux
fonctions de génération de contenu adaptatif.

\section{Étape 4 --- Niveau global pondéré}
\texttt{computeGlobal} agrège les niveaux de leçon en un niveau de chapitre,
pondéré par le nombre de réponses~:
\[
L_{\text{global}} = \frac{\sum_i L_i \times n_i}{\sum_i n_i}
\]
Elle identifie aussi la leçon la plus faible / la plus forte et signale une
\emph{lacune critique} dès qu'une leçon descend sous 30.

\section{Étape 5 --- Décroissance temporelle (oubli)}
\texttt{applyDecay} simule l'oubli en cas d'inactivité~:
\begin{itemize}
  \item $\leq 7$ jours : aucune décroissance.
  \item $> 7$ jours : facteur $= \max(0{,}10,\ 1 - 0{,}01\times(\text{jours}-7))$.
\end{itemize}
La décroissance est appliquée à l'ouverture d'une leçon et persistée en base.

\section{Enregistrement des réponses}
\begin{itemize}
  \item \filepath{src/hooks/useAdaptiveContent.ts} --- gère le contenu adaptatif,
        les compteurs de session, et appelle \texttt{recordAnswer} (qui applique
        \texttt{computeDelta} puis met à jour \texttt{student\_scores}).
  \item \filepath{src/lib/recordActivityAnswer.ts} --- fonction partagée permettant
        aux activités \emph{de chapitre} (non adaptatives) d'alimenter aussi le
        moteur de niveau, avec les mêmes signaux comportementaux.
  \item \filepath{src/lib/chatEvalTracker.ts} --- suit les évaluations issues du chat.
\end{itemize}

\begin{notebox}[\faLightbulb\ Comportement face à une mauvaise réponse]
Dans les exercices, une réponse fausse \textbf{n'affiche jamais la bonne réponse
automatiquement}~: l'interface se verrouille en rouge pendant 3 secondes, puis
l'élève peut réessayer. La solution détaillée n'apparaît que s'il clique
volontairement sur le bouton «~Voir la solution détaillée~» (ce qui compte alors
comme un abandon). Dans les quiz, l'élève sélectionne une option puis confirme
via le bouton «~Confirmer~».
\end{notebox}

\section{Synthèse du calcul du niveau}
\begin{enumerate}
  \item \textbf{Niveau initial} = score du test de positionnement.
  \item \textbf{Après chaque réponse} : $\Delta$ ELO (difficulté + temps) modulé
        par les signaux comportementaux $\rightarrow$ nouveau \texttt{current\_level}.
  \item \textbf{Niveau de chapitre} = moyenne pondérée des niveaux de leçon.
  \item \textbf{Score composite} = calibration de la difficulté du contenu IA.
  \item \textbf{Décroissance} = $-1\%$/jour au-delà de 7 jours d'inactivité.
\end{enumerate}

\chapter{Base de données Supabase}

La base est une instance \textbf{PostgreSQL} gérée par Supabase. Elle comporte
\textbf{21 tables} dans le schéma \texttt{public}, protégées par RLS. Cette
section décrit chaque table, ses colonnes principales et son rôle. Les
énumérations utilisées sont~:
\begin{itemize}
  \item \texttt{app\_role} : \texttt{student, parent, admin, pedago} ;
  \item \texttt{school\_level} : \texttt{5eme\_primaire, 1ere\_cem, 2eme\_cem,
        3eme\_cem, 4eme\_cem, premiere, seconde, terminale} ;
  \item \texttt{link\_status} : \texttt{pending, active, rejected}.
\end{itemize}

\begin{notebox}
Les relations entre tables sont \emph{logiques} (par \texttt{user\_id},
\texttt{chapter\_id}, \texttt{lesson\_id}\dots) et non matérialisées par des clés
étrangères strictes vers \texttt{auth.users} (recommandation Supabase). Le MCD du
chapitre~\ref{chap:mcd} les illustre.
\end{notebox}

% ---------- Macro pour décrire une table ----------
\newcommand{\dbtable}[2]{%
\subsection{\texttt{#1}}%
\emph{#2}%
}

\section{Utilisateurs et rôles}

\dbtable{profiles}{Profil applicatif de chaque utilisateur (extension de auth.users).}
\begin{longtable}{p{3.6cm}p{2.6cm}p{8cm}}
\toprule \textbf{Colonne} & \textbf{Type} & \textbf{Description} \\ \midrule \endhead
\texttt{id} & uuid & Identifiant (= id auth). \\
\texttt{email} & text & Adresse e-mail. \\
\texttt{email\_verified} & boolean & E-mail confirmé. \\
\texttt{first\_name / last\_name} & text & Nom et prénom. \\
\texttt{phone} & text & Téléphone. \\
\texttt{avatar\_url} & text & Photo de profil (bucket \texttt{avatars}). \\
\texttt{school\_level} & enum & Niveau scolaire. \\
\texttt{filiere} & text & Filière (secondaire). \\
\texttt{linking\_code} & text & Code de liaison parent--enfant. \\
\texttt{is\_active} & boolean & Compte actif. \\
\texttt{date\_of\_birth} & date & Date de naissance. \\
\texttt{wilaya / ville / ecole} & text & Localisation et établissement. \\
\bottomrule
\end{longtable}

\dbtable{user\_roles}{Rôles applicatifs (séparés du profil pour la sécurité).}
\begin{longtable}{p{3.6cm}p{2.6cm}p{8cm}}
\toprule \textbf{Colonne} & \textbf{Type} & \textbf{Description} \\ \midrule \endhead
\texttt{id} & uuid & Identifiant. \\
\texttt{user\_id} & uuid & Utilisateur concerné. \\
\texttt{role} & app\_role & Rôle (student/parent/admin/pedago). \\
\texttt{created\_at} & timestamptz & Date d'attribution. \\
\bottomrule
\end{longtable}

\dbtable{parent\_child\_links}{Liaison entre un parent et un enfant.}
\begin{longtable}{p{3.6cm}p{2.6cm}p{8cm}}
\toprule \textbf{Colonne} & \textbf{Type} & \textbf{Description} \\ \midrule \endhead
\texttt{parent\_id} & uuid & Parent. \\
\texttt{child\_id} & uuid & Enfant. \\
\texttt{status} & link\_status & État (pending/active/rejected). \\
\bottomrule
\end{longtable}

\dbtable{activity\_logs}{Journal d'audit des actions (via fonction \texttt{log\_activity}).}
\begin{longtable}{p{3.6cm}p{2.6cm}p{8cm}}
\toprule \textbf{Colonne} & \textbf{Type} & \textbf{Description} \\ \midrule \endhead
\texttt{user\_id} & uuid & Auteur de l'action. \\
\texttt{action} & text & Code de l'action. \\
\texttt{details} & jsonb & Détails contextuels. \\
\texttt{ip\_address} & text & Adresse IP (optionnelle). \\
\bottomrule
\end{longtable}

\section{Contenu pédagogique}

\dbtable{filieres}{Filières disponibles par niveau scolaire.}
\begin{longtable}{p{3.6cm}p{2.6cm}p{8cm}}
\toprule \textbf{Colonne} & \textbf{Type} & \textbf{Description} \\ \midrule \endhead
\texttt{school\_level} & enum & Niveau associé. \\
\texttt{code} & text & Code court (ex.\ \texttt{sciences}). \\
\texttt{name / name\_ar} & text & Libellé FR / AR. \\
\bottomrule
\end{longtable}

\dbtable{chapters}{Chapitres de cours, organisés par niveau, filière et matière.}
\begin{longtable}{p{3.6cm}p{2.6cm}p{8cm}}
\toprule \textbf{Colonne} & \textbf{Type} & \textbf{Description} \\ \midrule \endhead
\texttt{school\_level} & enum & Niveau scolaire. \\
\texttt{filiere\_id} & uuid & Filière (nullable). \\
\texttt{title / title\_ar} & text & Titre FR / AR. \\
\texttt{description} & text & Description. \\
\texttt{subject} & text & Matière (\texttt{math}). \\
\texttt{order\_index} & integer & Ordre d'affichage. \\
\bottomrule
\end{longtable}

\dbtable{lessons}{Leçons appartenant à un chapitre.}
\begin{longtable}{p{3.6cm}p{2.6cm}p{8cm}}
\toprule \textbf{Colonne} & \textbf{Type} & \textbf{Description} \\ \midrule \endhead
\texttt{chapter\_id} & uuid & Chapitre parent. \\
\texttt{title / title\_ar} & text & Titre FR / AR. \\
\texttt{content} & text & Contenu (Markdown/HTML). \\
\texttt{video\_url} & text & Vidéo associée. \\
\texttt{order\_index} & integer & Ordre. \\
\bottomrule
\end{longtable}

\dbtable{chapter\_exercises}{Exercices d'entraînement (réponse libre).}
\begin{longtable}{p{3.6cm}p{2.6cm}p{8cm}}
\toprule \textbf{Colonne} & \textbf{Type} & \textbf{Description} \\ \midrule \endhead
\texttt{chapter\_id / lesson\_id} & uuid & Rattachement. \\
\texttt{title} & text & Titre de l'exercice. \\
\texttt{statement} & text & Énoncé. \\
\texttt{expected\_answer} & text & Réponse attendue. \\
\texttt{accepted\_answers} & jsonb & Variantes acceptées. \\
\texttt{solution} & text & Solution détaillée. \\
\texttt{hint} & text & Indice. \\
\texttt{difficulty} & integer & Difficulté (1--5). \\
\bottomrule
\end{longtable}

\dbtable{chapter\_quizzes}{Questions de quiz (choix multiples).}
\begin{longtable}{p{3.6cm}p{2.6cm}p{8cm}}
\toprule \textbf{Colonne} & \textbf{Type} & \textbf{Description} \\ \midrule \endhead
\texttt{chapter\_id / lesson\_id} & uuid & Rattachement. \\
\texttt{question} & text & Énoncé de la question. \\
\texttt{options} & jsonb & Liste des choix. \\
\texttt{correct\_answer} & text & Bonne réponse. \\
\texttt{explanation} & text & Explication. \\
\texttt{hint} & text & Indice. \\
\texttt{difficulty} & integer & Difficulté (1--5). \\
\bottomrule
\end{longtable}

\dbtable{exams}{Examens / sujets par trimestre.}
\begin{longtable}{p{3.6cm}p{2.6cm}p{8cm}}
\toprule \textbf{Colonne} & \textbf{Type} & \textbf{Description} \\ \midrule \endhead
\texttt{school\_level} & enum & Niveau. \\
\texttt{trimester} & integer & Trimestre. \\
\texttt{title / title\_ar} & text & Titre. \\
\texttt{content} & jsonb & Contenu structuré. \\
\texttt{duration\_minutes} & integer & Durée. \\
\texttt{subject} & text & Matière. \\
\bottomrule
\end{longtable}

\section{Suivi de l'apprentissage}

\dbtable{student\_scores}{Cœur du moteur adaptatif : niveau et statistiques par élève/leçon/chapitre.}
\begin{longtable}{p{4.0cm}p{2.4cm}p{8cm}}
\toprule \textbf{Colonne} & \textbf{Type} & \textbf{Description} \\ \midrule \endhead
\texttt{user\_id} & uuid & Élève. \\
\texttt{chapter\_id / lesson\_id} & uuid & Portée (leçon, chapitre, ou globale si NULL). \\
\texttt{current\_level} & integer & Niveau courant (5--100). \\
\texttt{reading/quiz/exercise\_time\_seconds} & integer & Temps passé par type. \\
\texttt{total\_answers} & integer & Nombre total de réponses. \\
\texttt{correct\_answers} & integer & Réponses correctes. \\
\texttt{accuracy\_rate} & numeric & Taux de réussite (\%). \\
\texttt{streak} & integer & Série de bonnes réponses. \\
\texttt{assessment\_data} & jsonb & Drapeaux d'hésitation, données du test. \\
\texttt{periodic\_advice} & jsonb & Conseils périodiques. \\
\texttt{advice\_seen} & boolean & Conseil consulté. \\
\texttt{report\_first\_shown\_at} & timestamptz & Premier affichage du rapport. \\
\bottomrule
\end{longtable}

\dbtable{ai\_generated\_content}{Contenu (quiz/exercices/révision) généré par IA et calibré par niveau.}
\begin{longtable}{p{4.0cm}p{2.4cm}p{8cm}}
\toprule \textbf{Colonne} & \textbf{Type} & \textbf{Description} \\ \midrule \endhead
\texttt{user\_id} & uuid & Élève destinataire. \\
\texttt{chapter\_id / lesson\_id} & uuid & Rattachement. \\
\texttt{content\_type} & text & \texttt{quiz} / \texttt{exercise} / \texttt{revision}. \\
\texttt{content} & jsonb & Items générés. \\
\texttt{difficulty\_level} & integer & Niveau de calibration (composite). \\
\bottomrule
\end{longtable}

\dbtable{ai\_lesson\_comments}{Commentaires IA personnalisés après une session d'étude.}
\begin{longtable}{p{4.0cm}p{2.4cm}p{8cm}}
\toprule \textbf{Colonne} & \textbf{Type} & \textbf{Description} \\ \midrule \endhead
\texttt{user\_id} & uuid & Élève. \\
\texttt{lesson\_id / chapter\_id} & uuid & Leçon/chapitre concernés. \\
\texttt{level\_before / level\_after / level\_delta} & integer & Évolution du niveau. \\
\texttt{weak\_concepts / strong\_concepts} & jsonb & Lacunes et acquis détectés. \\
\texttt{message} & text & Commentaire (arabe, Markdown). \\
\texttt{link\_url} & text & Lien vers la leçon. \\
\bottomrule
\end{longtable}

\dbtable{parent\_reports}{Rapports de progression destinés aux parents.}
\begin{longtable}{p{4.0cm}p{2.4cm}p{8cm}}
\toprule \textbf{Colonne} & \textbf{Type} & \textbf{Description} \\ \midrule \endhead
\texttt{parent\_id / child\_id} & uuid & Parent et enfant. \\
\texttt{period\_start / period\_end} & timestamptz & Période couverte. \\
\texttt{report\_data} & jsonb & Données brutes. \\
\texttt{global\_success\_rate} & numeric & Taux de réussite global. \\
\texttt{global\_level} & integer & Niveau global. \\
\texttt{strong\_chapters / weak\_chapters} & jsonb & Forces et faiblesses. \\
\texttt{ai\_recommendations / summary} & text & Recommandations IA. \\
\texttt{report\_type} & text & Type de rapport. \\
\bottomrule
\end{longtable}

\section{Assistant IA (chat)}

\dbtable{chat\_conversations}{Historique des conversations avec l'assistant IA.}
\begin{longtable}{p{3.6cm}p{2.6cm}p{8cm}}
\toprule \textbf{Colonne} & \textbf{Type} & \textbf{Description} \\ \midrule \endhead
\texttt{user\_id} & uuid & Propriétaire. \\
\texttt{chapter\_id} & uuid & Contexte (nullable). \\
\texttt{title} & text & Titre de la conversation. \\
\texttt{messages} & jsonb & Fil de messages. \\
\bottomrule
\end{longtable}

\dbtable{chat\_usage}{Quotas d'utilisation quotidienne du chat (messages/images).}
\begin{longtable}{p{3.6cm}p{2.6cm}p{8cm}}
\toprule \textbf{Colonne} & \textbf{Type} & \textbf{Description} \\ \midrule \endhead
\texttt{user\_id} & uuid & Utilisateur. \\
\texttt{usage\_date} & date & Jour. \\
\texttt{message\_count} & integer & Nombre de messages. \\
\texttt{image\_count} & integer & Nombre d'images. \\
\bottomrule
\end{longtable}

\section{Abonnements et paiements}

\dbtable{subscription\_config}{Configuration des offres (prix individuel/famille).}
\begin{longtable}{p{3.6cm}p{2.6cm}p{8cm}}
\toprule \textbf{Colonne} & \textbf{Type} & \textbf{Description} \\ \midrule \endhead
\texttt{plan\_type} & text & Type d'offre. \\
\texttt{price\_single / price\_family} & integer & Tarifs. \\
\texttt{label} & text & Libellé. \\
\texttt{is\_active} & boolean & Offre active. \\
\bottomrule
\end{longtable}

\dbtable{subscription\_periods}{Périodes de validité des abonnements.}
\begin{longtable}{p{3.6cm}p{2.6cm}p{8cm}}
\toprule \textbf{Colonne} & \textbf{Type} & \textbf{Description} \\ \midrule \endhead
\texttt{label} & text & Libellé. \\
\texttt{start\_date / end\_date} & date & Bornes. \\
\texttt{is\_active} & boolean & Période active. \\
\bottomrule
\end{longtable}

\dbtable{student\_subscriptions}{Abonnement actif d'un élève (décompte de jours).}
\begin{longtable}{p{3.6cm}p{2.6cm}p{8cm}}
\toprule \textbf{Colonne} & \textbf{Type} & \textbf{Description} \\ \midrule \endhead
\texttt{user\_id} & uuid & Élève. \\
\texttt{activation\_code\_id} & uuid & Code utilisé. \\
\texttt{plan\_type} & text & Type d'offre. \\
\texttt{total\_days / days\_used} & integer/numeric & Crédit et consommation. \\
\texttt{is\_paused / paused\_at} & boolean/ts & Mise en pause. \\
\texttt{started\_at / last\_tick\_at} & timestamptz & Démarrage et dernier décompte. \\
\bottomrule
\end{longtable}

\dbtable{activation\_codes}{Codes prépayés d'activation d'abonnement.}
\begin{longtable}{p{3.6cm}p{2.6cm}p{8cm}}
\toprule \textbf{Colonne} & \textbf{Type} & \textbf{Description} \\ \midrule \endhead
\texttt{code} & text & Code unique. \\
\texttt{plan\_type} & text & Offre liée. \\
\texttt{is\_family} & boolean & Offre famille. \\
\texttt{status} & text & État (free/used\dots). \\
\texttt{created\_by / used\_by} & uuid & Créateur / utilisateur. \\
\texttt{payment\_id} & uuid & Paiement associé. \\
\bottomrule
\end{longtable}

\dbtable{payments}{Historique des paiements.}
\begin{longtable}{p{3.6cm}p{2.6cm}p{8cm}}
\toprule \textbf{Colonne} & \textbf{Type} & \textbf{Description} \\ \midrule \endhead
\texttt{user\_id} & uuid & Payeur. \\
\texttt{plan\_type / plan\_label} & text & Offre. \\
\texttt{amount} & integer & Montant. \\
\texttt{is\_family / children\_count} & bool/int & Offre famille. \\
\texttt{period\_id} & uuid & Période. \\
\texttt{status} & text & État du paiement. \\
\bottomrule
\end{longtable}

\section{Fonctions de base de données}
\begin{longtable}{p{4.8cm}p{9.6cm}}
\toprule \textbf{Fonction} & \textbf{Rôle} \\ \midrule \endhead
\texttt{has\_role(uid, role)} & Vérifie un rôle (SECURITY DEFINER, anti-récursion RLS). \\
\texttt{is\_parent\_of(parent, child)} & Vérifie une liaison parent--enfant active. \\
\texttt{user\_has\_any\_role(uid)} & Indique si l'utilisateur a au moins un rôle. \\
\texttt{handle\_new\_user()} & Trigger : crée profil + rôle à l'inscription. \\
\texttt{handle\_updated\_at()} & Trigger : met à jour \texttt{updated\_at}. \\
\texttt{log\_activity(uid, action, details)} & Insère une ligne d'audit. \\
\texttt{generate\_activation\_code()} & Génère un code d'activation aléatoire. \\
\texttt{check\_exercise\_answer(id, rep)} & Valide une réponse d'exercice (comparaison normalisée exacte). \\
\texttt{check\_quiz\_answer(id, rep)} & Valide une réponse de quiz. \\
\texttt{get\_student\_exercises / get\_student\_quizzes} & Renvoie les exercices/quiz d'un chapitre/leçon. \\
\bottomrule
\end{longtable}

\begin{infobox}[\faCheckCircle\ Validation exacte des réponses]
La fonction \texttt{check\_exercise\_answer} normalise la saisie (minuscules,
suppression des espaces) puis effectue une \textbf{comparaison exacte} avec la
réponse attendue et les variantes acceptées. Cela évite qu'une saisie partielle
(ex.\ «~1~») soit acceptée pour une réponse «~12~».
\end{infobox}

\chapter{Modèle conceptuel de données (MCD)}
\label{chap:mcd}

Ce chapitre présente le modèle conceptuel des données. Les associations sont
\emph{logiques}~: elles reposent sur des colonnes d'identifiant (\texttt{user\_id},
\texttt{chapter\_id}, \texttt{lesson\_id}\dots) plutôt que sur des contraintes de
clés étrangères strictes vers \texttt{auth.users} (bonne pratique Supabase).

\section{Diagramme entité–association}
Le schéma ci-dessous résume les principales entités et leurs cardinalités.

\begin{center}
\footnotesize
\begin{verbatim}
                          +---------------------+
                          |     auth.users      |
                          +----------+----------+
                                     | 1
              +----------------------+----------------------+
              | 1                    | 1                    | 1
        +-----v------+        +------v-------+       +-------v--------+
        |  profiles  |        |  user_roles  |       | student_scores |
        +-----+------+        +--------------+       +-------+--------+
              | 1                                            | *
              | (linking_code)                               |
        +-----v---------------+                              |
        | parent_child_links  |  * --- 1 (parent/child)      |
        +---------------------+                              |
                                                             |
   PROGRAMME (geré par pedago/admin)                         |
        +-----------+   1      *  +-----------+   1   *  +----v-----+
        | filieres  |----------->| chapters  |--------->| lessons  |
        +-----------+            +-----+-----+          +----+-----+
                                    | 1                    | 1
                       +------------+----------+   +-------+---------+
                       | 1                     |   | *               | *
                +------v-------+      +---------v---+        +--------v---------+
                | chapter_     |      | chapter_    |        | ai_generated_    |
                | quizzes      |      | exercises   |        | content          |
                +--------------+      +-------------+        +------------------+

   SUIVI & IA                         ABONNEMENTS & PAIEMENTS
   +---------------------+            +------------------------+
   | ai_lesson_comments  |            | subscription_config    |
   | chat_conversations  |            | subscription_periods   |
   | chat_usage          |            | student_subscriptions  |
   | parent_reports      |            | activation_codes       |
   | activity_logs       |            | payments               |
   +---------------------+            +------------------------+
              (rattachés à user_id)        (rattachés à user_id)
\end{verbatim}
\end{center}

\section{Cardinalités principales}
\begin{longtable}{p{4.5cm}p{1.6cm}p{4.5cm}p{2.8cm}}
\toprule
\textbf{Entité A} & \textbf{Card.} & \textbf{Entité B} & \textbf{Clé logique} \\
\midrule \endhead
profiles & 1 — 1 & auth.users & \texttt{id} \\
user\_roles & * — 1 & auth.users & \texttt{user\_id} \\
parent\_child\_links & * — 1 & profiles (parent) & \texttt{parent\_id} \\
parent\_child\_links & * — 1 & profiles (enfant) & \texttt{child\_id} \\
filieres & 1 — * & chapters & \texttt{filiere\_id} \\
chapters & 1 — * & lessons & \texttt{chapter\_id} \\
chapters & 1 — * & chapter\_exercises & \texttt{chapter\_id} \\
chapters & 1 — * & chapter\_quizzes & \texttt{chapter\_id} \\
lessons & 1 — * & chapter\_exercises & \texttt{lesson\_id} \\
lessons & 1 — * & chapter\_quizzes & \texttt{lesson\_id} \\
profiles (élève) & 1 — * & student\_scores & \texttt{user\_id} \\
lessons & 1 — * & student\_scores & \texttt{lesson\_id} \\
profiles (élève) & 1 — * & ai\_generated\_content & \texttt{user\_id} \\
profiles (élève) & 1 — * & ai\_lesson\_comments & \texttt{user\_id} \\
profiles & 1 — * & chat\_conversations & \texttt{user\_id} \\
parent + enfant & 1 — * & parent\_reports & \texttt{parent\_id, child\_id} \\
profiles (élève) & 1 — * & student\_subscriptions & \texttt{user\_id} \\
activation\_codes & 1 — 1 & student\_subscriptions & \texttt{activation\_code\_id} \\
profiles & 1 — * & payments & \texttt{user\_id} \\
\bottomrule
\end{longtable}

\section{Regroupement fonctionnel des entités}
\begin{itemize}
  \item \textbf{Identité \& accès}~: \texttt{profiles}, \texttt{user\_roles},
        \texttt{parent\_child\_links}, \texttt{activity\_logs}.
  \item \textbf{Programme}~: \texttt{filieres}, \texttt{chapters},
        \texttt{lessons}, \texttt{chapter\_exercises}, \texttt{chapter\_quizzes},
        \texttt{exams}.
  \item \textbf{Apprentissage adaptatif}~: \texttt{student\_scores},
        \texttt{ai\_generated\_content}, \texttt{ai\_lesson\_comments}.
  \item \textbf{Assistant IA}~: \texttt{chat\_conversations}, \texttt{chat\_usage}.
  \item \textbf{Suivi parental}~: \texttt{parent\_reports}.
  \item \textbf{Commerce}~: \texttt{subscription\_config},
        \texttt{subscription\_periods}, \texttt{student\_subscriptions},
        \texttt{activation\_codes}, \texttt{payments}.
\end{itemize}

\chapter{Fonctions serveur (Edge Functions IA)}

Toute la logique faisant appel à l'\textbf{intelligence artificielle} est
déportée dans des \textbf{Edge Functions Deno} hébergées par Supabase
(\filepath{supabase/functions/}). Ce choix garantit que les clés d'API ne sont
\textbf{jamais} exposées au navigateur~: elles sont lues côté serveur via
\texttt{Deno.env.get(...)}.

\section{Fournisseurs et modèles d'IA}
\begin{longtable}{p{5cm}p{9.5cm}}
\toprule \textbf{Fournisseur} & \textbf{Usage} \\ \midrule \endhead
\textbf{Lovable AI Gateway} & Passerelle principale vers Google Gemini
(\texttt{google/gemini-2.5-flash}, \texttt{google/gemini-3-flash-preview}).
Clé \texttt{LOVABLE\_API\_KEY}. \\
\textbf{Groq} & Modèle de repli rapide \texttt{llama-3.3-70b-versatile}.
Clé \texttt{GROQ\_API\_KEY}. \\
\textbf{Google Gemini (direct)} & Génération en masse via \texttt{GEMINI\_API\_KEY} /
\texttt{GEMINI\_API\_KEY\_2}. \\
\textbf{OpenRouter / Cloudflare AI} & Clés additionnelles disponibles
(\texttt{OPENROUTER\_API\_KEY}, \texttt{CLOUDFLARE\_AI\_API\_KEY}). \\
\bottomrule
\end{longtable}

\begin{notebox}[\faLock\ Stratégie de repli]
Plusieurs fonctions tentent d'abord Gemini (via la passerelle) puis basculent
automatiquement sur Groq (Llama) en cas d'erreur, afin de garantir la
disponibilité du service.
\end{notebox}

\section{Catalogue des fonctions}
\begin{longtable}{p{5cm}p{9.5cm}}
\toprule \textbf{Fonction} & \textbf{Rôle} \\ \midrule \endhead
\texttt{lovable-chat} & Assistant pédagogique IA (chatbot). Réponses
contextualisées et strictement encadrées par le programme de mathématiques.
Gemini avec repli Llama. \\
\texttt{generate-adaptive-content} & Génère quiz / exercices / révisions
calibrés au niveau composite de l'élève (Gemini 3 flash). \\
\texttt{generate-lesson-comment} & Produit un commentaire IA personnalisé en
arabe après une session d'étude (forces/lacunes, conseils). \\
\texttt{generate-remediation} & Génère des activités de remédiation ciblées
\emph{uniquement} sur les lacunes détectées de l'élève. \\
\texttt{generate-chapter-quizzes} & Génère des quiz pour tout un chapitre
(Gemini + repli Llama). \\
\texttt{generate-chapter-revision} & Génère des fiches de révision de chapitre. \\
\texttt{generate-parent-report} & Construit le rapport de progression destiné
aux parents (synthèse + recommandations IA). \\
\texttt{scheduled-parent-reports} & Tâche planifiée (pg\_cron) générant
quotidiennement les rapports pour chaque liaison parent--enfant active. \\
\texttt{enrich-chapter-lessons} & Enrichit le contenu des leçons d'un chapitre. \\
\texttt{bulk-generate-lesson-content} & Génère en masse 10 exercices + 10 quiz
pour une leçon et les insère en base. \\
\texttt{bulk-generate-all-terminale} & Génération par lot pour toutes les leçons
de Terminale Sciences. \\
\texttt{bulk-gen-terminale-gemini} & Variante de génération en masse via l'API
Gemini directe (\texttt{GEMINI\_API\_KEY\_2}). \\
\texttt{delete-user-account} & Suppression complète et sécurisée d'un compte
utilisateur (privilèges service\_role). \\
\bottomrule
\end{longtable}

\section{Secrets configurés}
Les secrets suivants sont disponibles côté Edge Functions (jamais exposés au
client)~:
\begin{itemize}
  \item IA~: \texttt{LOVABLE\_API\_KEY}, \texttt{GROQ\_API\_KEY},
        \texttt{GEMINI\_API\_KEY}, \texttt{GEMINI\_API\_KEY\_2},
        \texttt{OPENROUTER\_API\_KEY}, \texttt{CLOUDFLARE\_AI\_API\_KEY},
        \texttt{CLOUDFLARE\_ACCOUNT\_ID}.
  \item Supabase~: \texttt{SUPABASE\_URL}, \texttt{SUPABASE\_DB\_URL},
        \texttt{SUPABASE\_SERVICE\_ROLE\_KEY}, \texttt{SUPABASE\_ANON\_KEY},
        \texttt{SUPABASE\_PUBLISHABLE\_KEY}, \texttt{SUPABASE\_JWKS}\dots
\end{itemize}

\chapter{Sécurité et règles d'accès (RLS)}

\section{Modèle de sécurité}
La sécurité repose sur trois piliers~:
\begin{enumerate}
  \item \textbf{Authentification JWT} via Supabase Auth.
  \item \textbf{Rôles séparés} dans \texttt{user\_roles}, vérifiés par des
        fonctions \texttt{SECURITY DEFINER} (\texttt{has\_role},
        \texttt{is\_parent\_of}) --- évite l'élévation de privilèges et la
        récursion RLS.
  \item \textbf{Row-Level Security (RLS)} activée sur \emph{toutes} les tables du
        schéma \texttt{public}, complétée par des \texttt{GRANT} explicites.
\end{enumerate}

\section{Principes des politiques RLS}
\begin{longtable}{p{4.5cm}p{10cm}}
\toprule \textbf{Domaine} & \textbf{Règle d'accès} \\ \midrule \endhead
Données personnelles & Chaque utilisateur ne voit/modifie que ses propres lignes
(\texttt{auth.uid() = user\_id}) : \texttt{student\_scores},
\texttt{ai\_generated\_content}, \texttt{ai\_lesson\_comments},
\texttt{chat\_conversations}, \texttt{chat\_usage}. \\
Contenu pédagogique & Lecture réservée aux utilisateurs authentifiés ; gestion
(écriture) réservée aux \texttt{pedago} et \texttt{admin} : \texttt{chapters},
\texttt{lessons}, \texttt{chapter\_exercises}, \texttt{chapter\_quizzes},
\texttt{exams}, \texttt{filieres}. \\
Suivi parental & Un parent accède aux données de ses enfants \emph{uniquement}
si une liaison active existe (\texttt{is\_parent\_of}) :
\texttt{parent\_reports}, \texttt{ai\_lesson\_comments} (enfants). \\
Administration & Les administrateurs disposent de politiques \og manage\fg{}
(ALL) sur les codes, abonnements, journaux, etc. \\
Audit & \texttt{activity\_logs} : insertion par utilisateurs authentifiés,
lecture par le propriétaire et les administrateurs. \\
\bottomrule
\end{longtable}

\section{Bonnes pratiques appliquées}
\begin{itemize}
  \item \textbf{Aucun secret en base} : les clés d'API vivent dans les secrets
        Supabase, jamais dans les tables ni dans le client.
  \item \textbf{\texttt{service\_role} jamais exposé} au navigateur ; utilisé
        seulement à l'intérieur des Edge Functions.
  \item \textbf{Validation des réponses côté serveur} via des fonctions
        \texttt{SECURITY DEFINER} (\texttt{check\_exercise\_answer},
        \texttt{check\_quiz\_answer}) --- l'élève ne peut pas connaître la réponse
        en inspectant le client.
  \item \textbf{GRANT explicites} pour chaque table conformément aux exigences de
        PostgREST.
\end{itemize}

\begin{infobox}[\faExclamationTriangle\ Point de vigilance]
Les rôles ne doivent jamais migrer vers la table \texttt{profiles}. Toute
vérification d'autorisation côté client (\texttt{hasRole}) est un confort d'UX~;
la véritable barrière de sécurité reste la \textbf{RLS} et les fonctions
\texttt{SECURITY DEFINER} côté base.
\end{infobox}

\chapter{Déploiement et maintenance}

\section{Environnements}
\begin{longtable}{p{4cm}p{10.5cm}}
\toprule \textbf{Environnement} & \textbf{URL} \\ \midrule \endhead
Aperçu (preview) & \url{https://id-preview--d3d48394-d747-4de9-844c-541a6a3403d6.lovable.app} \\
Production & \url{https://academieplus.lovable.app} \\
\bottomrule
\end{longtable}

Le projet Supabase associé porte la référence \texttt{lfothlxoixayjiytwwqa}.

\section{Variables d'environnement}
Le fichier \filepath{.env} est rempli automatiquement avec~:
\begin{lstlisting}[language=bash]
VITE_SUPABASE_URL=...
VITE_SUPABASE_PUBLISHABLE_KEY=...   # cle publique (anon) - OK cote client
VITE_SUPABASE_PROJECT_ID=...
\end{lstlisting}
La clé publique (\texttt{anon}) est volontairement publique~; la protection
réelle des données provient de la RLS. Les clés \emph{privées} restent dans les
secrets Supabase.

\section{Cycle de déploiement}
\begin{itemize}
  \item \textbf{Frontend} : build Vite (\texttt{npm run build}) servi en statique.
  \item \textbf{Edge Functions} : déploiement \textbf{automatique} à chaque
        modification (aucune commande manuelle).
  \item \textbf{Migrations} : les fichiers SQL de \filepath{supabase/migrations/}
        sont versionnés et appliqués via l'outil de migration.
\end{itemize}

\section{Internationalisation}
L'application est bilingue \textbf{français / arabe} via \texttt{i18next}
(\filepath{src/i18n/}). Les libellés sont stockés dans
\filepath{src/i18n/locales/fr.json} et \filepath{src/i18n/locales/ar.json}.
L'interface élève est majoritairement en \textbf{arabe} (RTL), tandis que les
espaces d'administration et éditorial sont en français.

\section{Tests}
\begin{itemize}
  \item \textbf{Unitaires / composants} : Vitest + Testing Library
        (\filepath{src/test/}). Le moteur de niveau (\filepath{src/lib/levelEngine.ts})
        étant constitué de fonctions pures, il est idéal à couvrir par des tests.
  \item \textbf{Bout en bout} : Playwright (\filepath{playwright.config.ts}).
\end{itemize}

\section{Conseils pour les prochaines équipes}
\begin{enumerate}
  \item \textbf{Logique de niveau} : toute évolution du calcul du niveau doit
        passer par \filepath{src/lib/levelEngine.ts} (fonctions pures, testables)
        plutôt que d'être dispersée dans les composants.
  \item \textbf{Accès données} : privilégier les services (\filepath{src/services})
        et hooks (\filepath{src/hooks}) existants plutôt que des appels Supabase
        directs dans les composants.
  \item \textbf{Sécurité} : ne jamais déplacer les rôles dans \texttt{profiles}~;
        toujours créer les \texttt{GRANT} et politiques RLS dans la même migration
        qu'un \texttt{CREATE TABLE}.
  \item \textbf{IA} : ajouter toute nouvelle intégration IA dans une Edge Function,
        jamais côté client.
  \item \textbf{Nettoyage} : la racine du dépôt contient des scripts ponctuels
        historiques (\filepath{fix*.mjs}, \filepath{patch*.cjs}\dots) qui ne font
        pas partie du runtime et peuvent être archivés.
\end{enumerate}

\section*{Conclusion}
\addcontentsline{toc}{section}{Conclusion}
AcadémiePlus est une plateforme adaptative complète couvrant l'ensemble du
parcours d'un élève~: positionnement initial, cours, exercices/quiz calibrés,
assistant IA, suivi de progression et reporting parental. Son architecture
(React + Supabase + Edge Functions IA) sépare clairement présentation, logique
métier pure et accès données, ce qui en facilite l'évolution.


\end{document}
